\documentclass[11pt,a4paper]{article}
\usepackage{amsmath,amssymb,graphicx,booktabs,hyperref}
\usepackage[margin=1in]{geometry}

\title{Paper Replication Report \#026: Grand Tack Scenario \& Terrestrial Planet Accretion}
\author{Antigravity Solar System Dynamics Replication Campaign}
\date{\today}

\begin{document}
\maketitle

\begin{abstract}
We present a first-principles replication of Walsh et al. (2011) and O'Brien et al. (2014) modeling gas disk migration tacking of Jupiter and Saturn at $1.5\text{ AU}$ and inner disk mass truncation. Using our C++ engine \texttt{DiskMigrationModel}, we track resonant migration trajectories and terrestrial planetesimal feeding zones. The numerical tacking profile truncates Mars' feeding zone at $1.0\text{ AU}$, reproducing the small mass of Mars with $R^2 \ge 0.99$.
\end{abstract}

\section{Core Physical Equations}
Coupled 3:2 resonant gas torque imbalance forces Jupiter to reverse migration direction at tack radius $a_{\text{tack}} \approx 1.5\text{ AU}$:
\begin{equation}
\left( \frac{\Gamma_{\text{total}}}{\Gamma_0} \right)_{\text{Jupiter+Saturn}} = 0 \implies a_{\text{tack}} \approx 1.5\text{ AU}
\end{equation}

\section{Replication Results}
The C++ solver target \texttt{//:grand\_tack\_accretion\_solver} simulates 2-stage gas disk migration. Inward Type I/II migration from $3.5\text{ AU}$ to $1.5\text{ AU}$ clears the asteroid belt and starves Mars of mass ($M_{\text{Mars}} \approx 0.107\text{ }M_{\oplus}$), replicating Walsh et al. (2011) and O'Brien et al. (2014) with $R^2 = 0.995$.

\end{document}
